% !TITLE 长相思·汴水流
% !AUTHOR 白居易
% !DYNASTY 唐
% !GENRE 词
% !ORDER 51

\begin{poem}
汴水\textsuperscript{1}流，泗水\textsuperscript{2}流，流到瓜洲\textsuperscript{3}古渡头。吴山\textsuperscript{4}点点愁。
思悠悠\textsuperscript{5}，恨悠悠，恨到归时方始休。月明人倚楼。
\end{poem}

\begin{annotations}
\notetext{1}{汴水：古河流名，发源于河南，流经开封，是古代南北交通的重要水道。}
\notetext{2}{泗水：古河流名，发源于山东，流入淮河。汴水流入泗水，泗水汇入淮河，淮河再入长江，一路奔向江南。}
\notetext{3}{瓜洲：古渡口名，在今江苏扬州，位于长江边。瓜洲渡是从中原到江南的必经之路。}
\notetext{4}{吴山：泛指江南的山。江南古属吴地，所以称吴山。女子望着远方的吴山，山只是一点点模糊的轮廓，但那轮廓里全是愁。}
\notetext{5}{悠悠：悠长、绵延不绝的样子。思悠悠是思念没有尽头，恨悠悠是怨恨也没有尽头。两个叠词连用，写出了思念的绵长和厚重。}
\end{annotations}

\begin{appreciation}
这首词写的是等待。白居易写的，也是中国词史上把等待写得最干净的一首之一。

汴水流，泗水流，流到瓜洲古渡头——汴水、泗水都是北方的河。汴水汇进泗水，泗水汇进淮河，淮河再进长江，一路流到江南。三句连用三个流字，像河水推着河水，一刻不停。水的方向是明确的，一路向南；思念的方向也是明确的，跟着水走。瓜洲是渡口——渡口是那个人走掉的地方，也是他该回来的地方。

吴山点点愁——江南的山，远远看去只剩一点点轮廓，可那一点点全是愁。山不会愁，是看山的人愁，愁到把山都看成了愁。

思悠悠，恨悠悠，恨到归时方始休——恨不是讨厌，是想你想得有点怨你的那个恨。可这恨是有尽头的，尽头就是他回来的那天。恨到归时方始休，说破了等待的人的心事：我恨你，是因为我相信你会回来。

月明人倚楼——不写了。月光，楼，一个人。等待的人不需要说词，一个姿势就够了。

还记得《行行重行行》吗？那里的女子等得“衣带日已缓”，这里的女子等得月光都白了。两千年前和一千年前，等待的姿势一模一样。这本书后面还有一首《长相思》，纳兰性德写的，人在关外，风雪夜里想家，说“故园无此声”——同一个词牌，一边站着想家的人，一边站着等人的人，中间隔着整个天下。

等你将来去很远的地方读书，家里就会有人月明人倚楼了。隔三差五报个平安，别让她等。
\end{appreciation}
